As mentioned in \cref{sec:background}\todo{discuss this in \cref{sec:background}}, origami robots—and by extension, soft robots—are a suitable use-case for lifelong learning techniques.
To evaluate the method described in \cref{sec:methodology}, a robot was designed that serves as a good analogue to the toy problem, whilst also using kresling origami as a key part of its functionality.
\Cref{fig:robot-overview} provides important nomenclature for reference.

\begin{figure}
    \newcounter{annotations}
    \centering
    \begin{subfigure}[t]{\linewidth}
        \centering
        \missingfigure[figwidth=0.9\linewidth]{Outside view of the robot with nomenclature labels}
        \begin{multicols}{2}
            \small
            \begin{enumerate}
                \item KOS               \label{item:kos}
                \item Actuation module  \label{item:actuation}
                \item Actuation arm     \label{item:arm}
                \item Connecting link   \label{item:link}
                \item Eye screw         \label{item:eye}
                \item Arduino           \label{item:arduino}
                \setcounter{annotations}{\value{enumi}}
            \end{enumerate}
        \end{multicols}
        \caption{The Kresling origami robot.}
        \label{fig:nomen-robot}
        \vspace*{1em}
    \end{subfigure}
    \begin{subfigure}[b]{\linewidth}
        \centering
        \missingfigure[figwidth=0.9\linewidth]{Inside view of the robot with nomenclature labels}
        \begin{multicols}{2}
            \small
            \begin{enumerate}
                \setcounter{enumi}{\value{annotations}}
                \item Servo motor   \label{item:servo}
                \item Servo arm     \label{item:servo-arm}
                \item Metal spring  \label{item:spring}
                \item Actuation arm
                \item Angle sensor  \label{item:sensor}
                \item Magnet        \label{item:magnet}
            \end{enumerate}
        \end{multicols}
        \caption{Actuation mechanism.}
        \label{fig:nomen-actuation}
    \end{subfigure}
    \caption{Annotated views of the robot and its most important actuation mechanisms}
    \label{fig:robot-overview}
\end{figure}

The robot is a simple crawler consisting of two Kresling origami springs (KOS) (\cref{fig:nomen-robot},~\cref{item:kos}) that are mirror images of each other, connected in series, with the actuation module (\ref{item:actuation}) in between.
Both springs can expand and contract in sync, producing a crawling-like motion when combined with specially designed ``feet''.
To demonstrate target adaptation, however, the movement of the robot as a whole is not very important.
Rather, the state of the KOSs is the state that will be controlled.

A KOS can be actuated both by linear motion (pushing or pulling) and by rotational motion (twisting).
Since the two motions are directly related (\cref{subsec:kresling-origami})\todo{\cref{subsec:kresling-origami}}, a contracting or expanding motion can be measured as the angle of the KOS' twist.
The twist is induced by turning a rotating arm (\ref{item:arm}) which is connected to the ends of the KOSs by two links (\ref{item:link}) that hook into eye-screws (\ref{item:eye}).
To turn the arm, a micro servo motor (\ref{item:servo}) is employed.
In the toy problem, the controller yields a force to exert on the system in order to reach the target position.
However, most micro servo's come with a near perfect position controller built-in, which cannot be accessed or changed.
Therefore, to bridge the gap between the simulation and the real-world implementation, the position output of the servo is converted to a force output by the use of a small, metal spring (\ref{item:spring}).
The servo thus does not directly act on the actuation arm, but rotates a servo arm (\ref{item:servo-arm}) instead, which pulls the spring and thereby exerts a force on the actuation arm.
As previously mentioned, the state of the KOS can be measured as the angle of twist, since it is directly related to its contraction.
Likewise, assuming perfectly rigid links with negligible clearance between the hooks and eye-screws, the angle of twist is directly related to the angle of the actuation arm.
This angle is measured using a contactless sensor (\ref{item:sensor}) that can measure the on-axis rotation of a small magnet (\ref{item:magnet}) attached to the actuation arm.


